% ============================================================
\section{Stakeholders (STK)}
% ============================================================

\subsection{Criterio de inclusión}

\begin{itemize}
  \item C1. Tiene intereses propios. Un sistema o servicio no tiene intereses; corresponde a una dependencia (DEP).
  \item C2. Su interés puede modificar una decisión de arquitectura.
  \item C3. No es redundante respecto de otro stakeholder ya identificado.
\end{itemize}

\begin{tabla}{|L{3.9cm}|L{2.75cm}|Y|}
\hline
\filacab \cab{Elemento} & \cab{Decisión} & \cab{Justificación} \\ \hline
\endhead
Jugadores adultos & Fusionado en STK-01 &
Falla C3. Usa las mismas pantallas, permisos y protocolo que el resto de jugadores. Su interés por la profundidad estratégica es un concern del mismo stakeholder, no un stakeholder distinto. \\ \hline
Equipo de desarrollo (genérico) & Eliminado &
Falla C3. Es la suma de roles ya listados por separado; conservarlo duplica intereses. \\ \hline
Diseñador de sonido y música & Fusionado en STK-08 &
Falla C3. Comparte los concerns del equipo gráfico: recursos propios o con licencia y localización de contenido. \\ \hline
Proveedor de servidores & Movido a DEP-04 &
Falla C1. No tiene interés en la arquitectura; el sistema depende de su disponibilidad. \\ \hline
Aplicación autenticadora, correo, SQL Server & Movidos a DEP &
Falla C1. Son sistemas de los que se depende, no partes interesadas. \\ \hline
Arquitecto de software & Incorporado como STK-05 &
Es el responsable de resolver las tensiones identificadas y de justificar las decisiones estructurales. \\ \hline
Auditor del proyecto & Incorporado como STK-16 &
La restricción de incrementos auditables semanales implica un actor que exige trazabilidad y evidencia. \\ \hline
Equipo de operaciones y despliegue & Incorporado como STK-14 &
El sistema se despliega en línea y se entrega en incrementos semanales; requiere un responsable de disponibilidad y despliegue. \\ \hline
\end{tabla}

% ------------------------------------------------------------
\subsection{Usuarios y clientes}

\begin{tabla}{|L{1.05cm}|L{3.0cm}|L{4.7cm}|Y|}
\hline
\filacab \cab{ID} & \cab{Stakeholder} & \cab{Interés principal} & \cab{Por qué influye en la arquitectura} \\ \hline
\endhead
STK-01 & Jugador (8 años o más) &
Reglas fáciles de aprender, controles accesibles, partidas fluidas y profundidad estratégica suficiente. &
Fija el piso de usabilidad (un niño de 8 años debe entender por qué una jugada es inválida) y el techo de latencia. Obliga a validación local inmediata más confirmación del servidor. \\ \hline
STK-02 & Padre o tutor &
Contenido seguro, protección del menor y control sobre lo que hace su hijo. &
Puede vetar el uso del juego. Impone filtro de chat, moderación y recolección mínima de datos personales, lo que afecta el modelo de cuentas y el diseño del chat. \\ \hline
\end{tabla}

% ------------------------------------------------------------
\subsection{Negocio y dirección}

\begin{tabla}{|L{1.05cm}|L{3.0cm}|L{4.7cm}|Y|}
\hline
\filacab \cab{ID} & \cab{Stakeholder} & \cab{Interés principal} & \cab{Por qué influye en la arquitectura} \\ \hline
\endhead
STK-03 & Profesor o cliente &
Cumplimiento de objetivos y entrega puntual. &
Fija el límite de 14 semanas y las auditorías semanales, lo que obliga a una arquitectura incremental y a priorizar el núcleo jugable sobre lo opcional. \\ \hline
STK-04 & Director del proyecto &
Cumplimiento de alcance, presupuesto y calendario. &
Es quien recorta alcance. Presiona hacia soluciones simples y decide si las funcionalidades adicionales, como la IA y la conexión LAN, entran o salen. \\ \hline
\end{tabla}

% ------------------------------------------------------------
\subsection{Desarrollo y arquitectura}

\begin{tabla}{|L{1.05cm}|L{3.0cm}|L{4.7cm}|Y|}
\hline
\filacab \cab{ID} & \cab{Stakeholder} & \cab{Interés principal} & \cab{Por qué influye en la arquitectura} \\ \hline
\endhead
STK-05 & Arquitecto de software &
Estructura coherente, modificabilidad y decisiones justificadas. &
Decide la separación cliente-servidor, el diseño del protocolo propio sobre TCP y quién tiene autoridad sobre el estado de la partida. \\ \hline
STK-06 & Diseñador del videojuego &
Mecánicas originales y equilibradas. &
Las reglas cambiarán durante el balanceo, lo que favorece un motor de reglas parametrizable. Además, la originalidad es una exigencia legal y no solo creativa. \\ \hline
STK-07 & Programadores &
Requisitos claros, recursos y tiempo suficientes. &
Los requisitos aún ambiguos, como la existencia de IA o de conexión LAN, obligan a diseñar puntos de extensión en lugar de comprometerse con una solución única. \\ \hline
STK-08 & Equipo de arte, interfaz y sonido &
Identidad visual y sonora propia e interfaz fácil de usar. &
La interfaz debe soportar textos de longitud variable entre inglés y español, lo que prohíbe los layouts de ancho fijo. La identidad propia viene impuesta por la propiedad intelectual de Quoridor. \\ \hline
STK-09 & Diseñador de experiencia de usuario &
Facilidad de aprendizaje y buena experiencia de uso. &
Exige tutorial, retroalimentación inmediata de jugadas ilegales y previsualización de barreras, por lo que el cliente necesita validar reglas localmente y no solo el servidor. \\ \hline
STK-10 & Equipo de pruebas (QA) &
Estabilidad, ausencia de errores y equilibrio. &
Necesita ejecutar y reproducir partidas sin interfaz, lo que obliga a separar el motor de reglas de la interfaz y a que las partidas sean deterministas y registrables. \\ \hline
STK-11 & Equipo de localización &
Traducciones claras y culturalmente apropiadas. &
Obliga a Unicode, a externalizar todas las cadenas fuera del código y a formatear fechas y números por región. Es transversal y no puede agregarse al final. \\ \hline
\end{tabla}

% ------------------------------------------------------------
\subsection{Operación y soporte}

\begin{tabla}{|L{1.85cm}|L{2.9cm}|L{4.45cm}|Y|}
\hline
\filacab \cab{ID} & \cab{Stakeholder} & \cab{Interés principal} & \cab{Por qué influye en la arquitectura} \\ \hline
\endhead
STK-12 \textit{(condicionado a Q-06)} & Equipo de soporte y moderación &
Atender reportes y sancionar conductas inapropiadas. &
Solo existe si se opta por moderación con revisión humana. En ese caso requiere historial de chat, identidad de cuenta trazable y un módulo administrativo. \\ \hline
STK-13 & Administrador de base de datos &
Integridad, disponibilidad y recuperación de los datos. &
Define el esquema en SQL Server, los respaldos y la recuperación. Su interés compite con la latencia: escribir cada jugada da trazabilidad pero cuesta tiempo. \\ \hline
STK-14 & Equipo de operaciones y despliegue &
Que el sistema sea desplegable y esté disponible. &
El modo en línea exige un servidor disponible de forma continua y entregas semanales, lo que empuja hacia configuración externa y despliegue automatizado. \\ \hline
\end{tabla}

% ------------------------------------------------------------
\subsection{Seguridad, auditoría y regulación}

\begin{tabla}{|L{1.05cm}|L{3.0cm}|L{4.7cm}|Y|}
\hline
\filacab \cab{ID} & \cab{Stakeholder} & \cab{Interés principal} & \cab{Por qué influye en la arquitectura} \\ \hline
\endhead
STK-15 & Especialista en ciberseguridad &
Seguridad de las cuentas y prevención de trampas. &
Exige servidor autoritativo (el cliente no puede decidir la validez de una jugada), protección del tráfico sobre TCP y segundo factor. El servidor autoritativo es una decisión estructural mayor. \\ \hline
STK-16 & Auditor del proyecto &
Trazabilidad y evidencia verificable cada semana. &
Obliga a que cada incremento sea funcional y demostrable y a mantener viva la trazabilidad entre stakeholders, concerns y decisiones, lo que condiciona el orden de construcción de los módulos. \\ \hline
STK-17 & Asesor legal y de propiedad intelectual &
Uso legal de nombres, diseños, reglas y recursos. &
Puede obligar a cambiar nombre, arte o redacción de reglas en cualquier momento, por lo que conviene aislar marca y textos del código. También regula el tratamiento de datos de menores. \\ \hline
\end{tabla}

% ------------------------------------------------------------
\subsection{Externos con capacidad de decisión}

\begin{tabla}{|L{1.45cm}|L{2.9cm}|L{4.7cm}|Y|}
\hline
\filacab \cab{ID} & \cab{Stakeholder} & \cab{Interés principal} & \cab{Por qué influye en la arquitectura} \\ \hline
\endhead
STK-18 & Gigamic y Mirko Marchesi &
Protección de su marca y propiedad intelectual. &
Nunca usará el sistema, pero tiene poder legal para exigir cambios o el cese del proyecto. Por eso es un stakeholder y no dependencia. \\ \hline
STK-19 (restricción legal) & Plataformas de distribución &
Cumplimiento de normas técnicas, legales y de contenido. &
Imponen clasificación por edad, requisitos de empaquetado y políticas de chat y de datos de menores, lo que condiciona el formato del entregable. \\ \hline
\end{tabla}
